\documentclass[11pt]{article}
\usepackage[margin=1in]{geometry}
\usepackage{amsmath,amssymb,amsthm,mathtools}
\usepackage[T1]{fontenc}
\usepackage{lmodern}
\usepackage[hidelinks]{hyperref}
\newtheorem{theorem}{Theorem}
\newtheorem{lemma}[theorem]{Lemma}
\newtheorem{proposition}[theorem]{Proposition}
\theoremstyle{remark}\newtheorem{remark}[theorem]{Remark}
\newcommand{\R}{\mathbb R}
\newcommand{\E}{\mathbb E}
\newcommand{\Prob}{\mathbb P}
\newcommand{\dd}{\,dx}
\newcommand{\Div}{\operatorname{div}}
\newcommand{\vol}{\operatorname{vol}}
\newcommand{\op}{\mathrm{op}}
\newcommand{\HS}{\mathrm{HS}}
\newcommand{\Id}{\mathrm I}
\newcommand{\T}{\mathsf T}
\newcommand{\tr}{\mathsf T}
\newcommand{\norm}[1]{\left\lVert#1\right\rVert}
\title{Consistent spectral recovery of smooth reversible\newline diffusion tensors at a fixed lag}
\author{Alec Kriebel\\Independent researcher\\
\small\href{https://orcid.org/0009-0001-9320-500X}{ORCID: 0009-0001-9320-500X}}
\date{October 5, 2026}
\begin{document}
\maketitle
\begin{abstract}
We give an almost surely consistent spectral-equation estimator for a general
smooth symmetric diffusion tensor and an unknown smooth invariant density on
a known bounded connected smooth domain. Observations are stationary exact
positions at one known fixed positive lag; reflection is conormal, and fixed
ellipticity and density bounds are known. A symmetric finite-rank estimate of
the stationary pair kernel supplies all its positive empirical eigenpairs.
Transition-eigenvalue fourth powers and vanishing ridge regularization give a
basis-invariant fit. We prove full spectral-jet excitation and construct the
required derivative-consistent kernels from the dependent data. The tensor
and its separately fitted divergence converge locally uniformly; an
ellipticity-clipped tensor converges globally in $L^2$. This provides a
convergence analysis for the spectral reconstruction question posed by Rei\ss{}
in OWR24/2017, in the explicit smooth model stated here. No rate or numerical
implementation is claimed. Spectral fitting and eigenvalue-power weighting
have established predecessors.
\end{abstract}

\section{Problem, model and contribution}
Rei\ss{}'s contribution with Chorowski, Gobet and Hoffmann asks for convergence
analysis of reconstruction of a matrix coefficient from empirical generator
eigenpairs and invariant density \cite[pp.~1507--1509]{reiss}. The report
leaves the estimator and retained spectral count to be specified. We provide
both, for the following class.

Let $D\subset\R^d$, $d\geq2$, be known, bounded and connected with
$C^\infty$ boundary. Fix a known lag $\Delta>0$. The unknown fields
$S=S^\tr$ and $\mu$ are $C^\infty(\overline D)$ and satisfy known bounds
\begin{equation}\label{bounds}
 0<\ell\Id\leq S(x)\leq\Lambda\Id,\qquad
 0<c\leq\mu(x)\leq C,\qquad \int_D\mu\dd=1.
\end{equation}
Here smoothness on the closed domain means smoothness up to its boundary.
The stationary reversible reflected diffusion has generator
\begin{equation}\label{generator}
 Lu=\mu^{-1}\Div(S\nabla u),\qquad n\cdot S\nabla u=0
 \quad\hbox{on }\partial D.
\end{equation}
We observe $Y_i=X_{i\Delta}$, $0\leq i\leq N$, exactly, without observations
between these times. The covariance in the interior stochastic differential
equation is $\Sigma=2S/\mu$ and its drift is
$\beta=(\Div S)/\mu$, where $(\Div S)_j=\sum_i\partial_iS_{ij}$.
The boundary condition in \eqref{generator} is essential: ordinary normal
reflection cannot be substituted for anisotropic conormal reflection.

Let $p_\Delta$ be the transition density relative to Lebesgue measure and
$q(x,y)=\mu(x)p_\Delta(x,y)$ the symmetric stationary pair density.
The estimator below fits the eigenfunction equations
\begin{equation}\label{pde}
 \Div(S\nabla u_k)=\nu_k\mu u_k,
 \qquad \kappa_k=e^{\Delta\nu_k}>0.
\end{equation}
Its contribution is consistency of this specified empirical spectral fit,
including an actual statistical construction and a proved excitation
condition. It is pointwise in every fixed $(S,\mu)$ satisfying
\eqref{bounds}, with no known smoothness norm or known spectral gap required
for its tuning. We do not establish rough-coefficient or nonreversible
results, additional sensor-noise inference, a convergence rate, minimax
optimality, computational efficiency, or error control for approximate
quadrature and eigensolvers.

\paragraph{Established foundations and priority scope.}
Fixed-lag identification of a reversible generator is established by
Hansen--Scheinkman \cite[Proposition~5.5]{hs}; multivariate spectral and
Dirichlet-form foundations appear in Chen--Hansen--Scheinkman \cite{chs}.
Crommelin--Vanden-Eijnden already proposed nonparametric drift/tensor spectral
residual fitting \cite{cv2006}; their 2011 author manuscript uses powers of
transition eigenvalues as weights \cite[Eq.~(4.7)]{cv2011}. Neither idea is
introduced here. Gobet--Hoffmann--Rei\ss{} and Chorowski develop scalar
one-dimensional spectral inference \cite{ghr,chorowski}. Nickl,
Giordano--Wang and Alberti--Barnes--Jambhale--Nickl supply multidimensional
scalar-conductivity fixed-lag theory \cite{nickl,gw,abjn}. In particular,
\cite[Lemma~3.8]{abjn} is a close gradient-excitation predecessor for our
cutoff-jet argument and credits earlier work of Alberti. Deterministic
matrix inversion from internal solutions has gradient/Hessian and
variational predecessors \cite{bal,kohn}.

A dated bounded primary-source audit located no theorem encompassing the
particular tensor spectral-consistency result proved here. This is not a
worldwide priority guarantee. In particular the full 2011 text inspected was
the CWI author manuscript revised August 31, 2011; the publisher-final text
was inaccessible, and identity with it is not assumed. Statements about its
deferred sample/discretization and infinite-expansion analyses refer to that
manuscript \cite[Sections~3.2 and~4.2]{cv2011}. Further edition and access
limits are recorded in the supplement. Classical nonlinear inverse
regularization \cite{lenzen} also yields an alternative estimator-existence
argument when its diffusion forward-map hypotheses are supplied; a separately
verified derivation of that application accompanies this note. That current
derivation is not a located earlier publication and does not establish
consistency of the spectral fit studied here. No claim of being the first
possible consistent tensor estimator is made.

\section{The empirical estimator}
We first specify a deterministic smoothing kernel that depends only on $D$.
The notation $C^{a,b}$ denotes the supremum of all mixed derivatives of
orders at most $a$ and $b$ in the respective variables.

\begin{lemma}[Known-domain smoothing]\label{kernel}
There are real, jointly Borel kernels $K_h(x,z)$, $z\in D$, with $x$ in a
fixed neighborhood of $\overline D$, smooth in $x$, such that
$R_h f(x)=\int_DK_h(x,z)f(z)\,dz$ converges to $f$ in
$C^2(\overline D)$ for $f\in C^3(\overline D)$ as $h\downarrow0$.
For each fixed derivative order $r$,
\begin{equation}\label{kernelbound}
 \sup_{x,z}|\partial_x^\alpha K_h(x,z)|\leq A_rh^{-d-|\alpha|}
 \quad (|\alpha|\leq r).
\end{equation}
On smooth functions of two variables, $R_h\otimes R_h$ converges to the
identity in $C^{2,2}(\overline D\times\overline D)$.
\end{lemma}
\begin{proof}
Choose a finite smooth boundary atlas and partition of unity with chart
cutoffs. In a flattened chart with inward coordinate $t\geq0$, extend a
cutoff function to a small exterior collar by
\begin{equation}\label{extension}
 (Ef)(x',t)=\sum_{a=1}^7\alpha_a f(x',-at),\quad t<0,
 \qquad (\alpha_1,\ldots,\alpha_7)=(28,-112,210,-224,140,-48,7).
\end{equation}
The identities $\sum_{a=1}^7\alpha_a(-a)^k=1$ for $0\leq k\leq6$
match derivatives through order six. Take the collar small enough for all
reflected arguments to remain in their charts, and add exterior cutoffs.
Interior pieces extend by zero after their cutoff. Summing gives a linear
compactly supported extension equal to $f$ on $D$ and bounded in each
$C^r$ norm, $r\leq6$.

Let $\eta$ be a fixed real $C_c^\infty(\R^d)$ mollifier of integral one
and set $R_h f=\eta_h*Ef$ on the indicated neighborhood. In each reflected
integral change variables back to $z\in D$. The chart maps have bounded
Jacobians bounded away from zero on their cutoff supports; the resulting
finite sum is a jointly Borel kernel. Differentiating the mollifier gives
\eqref{kernelbound} for any fixed $r$. Mollification of a $C^3$ extension
converges in $C^2$. Apply the extension independently in both variables to
obtain the mixed assertion. The kernels can be signed and need not preserve
mass; neither positivity nor exact mass preservation is used.
\end{proof}

Set $K_h(x,z)=0$ for $z\in\partial D$. This is a Borel convention
for every sample in $\overline D$; under stationarity,
$\Prob(Y_i\in\partial D)=0$ for every $i$, so it changes no
almost-sure conclusion.

For $N\geq2$ put $j=\lfloor\log_2N\rfloor$, $n=2^j$ and
$h=h_0/(j+1)$, where $h_0>0$ is a sufficiently small constant chosen from
the known charts. Use the first $n+1$ states and define
\begin{align}
 \widetilde\mu_j(x)&=\frac1{n+1}\sum_{i=0}^nK_h(x,Y_i),\label{density}\\
 q_j(x,y)&=\frac1{2n}\sum_{i=0}^{n-1}
 \big[K_h(x,Y_i)K_h(y,Y_{i+1})+
       K_h(y,Y_i)K_h(x,Y_{i+1})\big].\label{pair}
\end{align}
Choose a fixed smooth saturation $\psi:\R\to[c/2,2C]$ equal to the
identity on $[3c/4,3C/2]$, and set
\begin{equation}\label{saturation}
 \mu_j=\frac{\psi(\widetilde\mu_j)}{Z_j},\qquad
 Z_j=\int_D\psi(\widetilde\mu_j)\dd.
\end{equation}
In every sample, $\mu_j\geq c/(4C\vol(D))$. It is smooth, positive and
normalized. The kernel $q_j$ is real, symmetric and smooth, of rank at most
$n+1$, although it can be signed, indefinite, and have a marginal different
from $\mu_j$.

On $H=L^2(D,dx)$ define the real self-adjoint finite-rank operator $T_j$
with kernel $q_j(x,y)/\sqrt{\mu_j(x)\mu_j(y)}$. For every positive eigenvalue
$\kappa_{j,k}$ choose a real orthonormal basis $\phi_{j,k}$ of its
eigenspace and put
\begin{equation}\label{empiricalmodes}
 u_{j,k}=\phi_{j,k}/\sqrt{\mu_j},\qquad
 \nu_{j,k}=\Delta^{-1}\log\kappa_{j,k}.
\end{equation}
All positive eigenvalues are retained; negative and zero ones are omitted.
Nonzero modes are smooth because they lie in the span of the smooth kernel
sections. No labeling relative to population eigenvectors is imposed.

For a smooth $u$ define the column jet and matching unknown vector
\begin{align}
 Ju&=\big((\partial_{ii}u)_i,(2\partial_{ij}u)_{i<j},
                        (\partial_i u)_i\big)^\tr,\label{jet}\\
 \theta&=\big((S_{ii})_i,(S_{ij})_{i<j},(\Div S)_i\big)^\tr
        \in\R^m,\qquad m=d(d+1)/2+d.\notag
\end{align}
Thus $(Ju)^\tr\theta=\Div(S\nabla u)$. For each $x\in D$ let
\begin{equation}\label{fit}
 \widehat\theta_j(x)=\underset{z\in\R^m}{\arg\min}
 \left\{\sum_{\kappa_{j,k}>0}\kappa_{j,k}^4
 \big[(Ju_{j,k}(x))^\tr z-\nu_{j,k}\mu_j(x)u_{j,k}(x)\big]^2
       +\lambda_j|z|^2\right\},\qquad \lambda_j=(j+1)^{-1}.
\end{equation}
The minimizer is unique even if there are no positive modes. Write its tensor
part as the symmetric matrix $\widehat S_j$, and its remaining part as
$\widehat v_j$. The latter estimates $\Div S$ separately; it is not the
divergence of $\widehat S_j$. Let $\widehat S_j^c$ be the pointwise
Frobenius projection of $\widehat S_j$ onto
$\{M=M^\tr:\ell\Id\leq M\leq\Lambda\Id\}$. At sample size $N$, use
these estimators indexed by $j$; any fixed admissible tensor may be used for
the finitely many smaller sample sizes.

\begin{theorem}[Fixed-lag tensor consistency]\label{main}
Under \eqref{bounds}--\eqref{generator} and the stationary exact observation
model above, the specified estimators are measurable and, on one
probability-one event, for every compact $E\Subset D$,
\begin{equation}\label{local}
 \norm{\widehat S_j-S}_{C(E)}+
 \norm{\widehat v_j-\Div S}_{C(E)}\longrightarrow0,
 \qquad
 \norm{\widehat S_j^c-S}_{L^2(D)}\longrightarrow0.
\end{equation}
Also $\mu_j\to\mu$ in $C^2(\overline D)$, so
$2\widehat S_j/\mu_j$ and $\widehat v_j/\mu_j$ recover the covariance and
drift locally uniformly. The limits hold along all $N\to\infty$ by the
dyadic prefix convention. Any positive deterministic $\lambda_j\to0$ in
\eqref{fit} gives the same consistency conclusion.
\end{theorem}

\section{Derivative convergence from the dependent observations}
We give the statistical part explicitly, including overlapping pairs.
The form of $-L$ on $L^2(\mu\dd)$ is
$\mathcal E(f,f)=\int_D\nabla f^\tr S\nabla f\dd$. If
$\int_D f\mu\dd=0$, the Neumann Poincaré constant $C_D$ gives
\[
 \int_D f^2\mu\dd
 =\min_a\int_D(f-a)^2\mu\dd
 \leq C\min_a\int_D(f-a)^2\dd
 \leq CC_D\int_D|\nabla f|^2\dd.
\]
Consequently $P=e^{\Delta L}$ contracts mean-zero functions by
$\rho\leq\exp[-\Delta\ell/(CC_D)]<1$.

For a bounded pair observable $F(Y_i,Y_{i+1})$, center it as $Z_i$.
Let $r(z)=\E[Z_0\mid Y_1=z]$ and $s(z)=\E[Z_0\mid Y_0=z]$.
They have mean zero and $L^2(\mu)$ norms at most $\sqrt{\operatorname{Var}F}$.
Stationarity and the Markov property, conditioning at the endpoints, give
\begin{equation}\label{covariance}
 \operatorname{Cov}(Z_0,Z_a)=\langle r,P^{a-1}s\rangle_{L^2(\mu)},
 \qquad a\geq1.
\end{equation}
This includes the overlapping case $a=1$. Hence
$|\operatorname{Cov}(Z_0,Z_a)|\leq\rho^{a-1}\operatorname{Var}F$, and
\begin{equation}\label{variance}
 \operatorname{Var}\left(\frac1n\sum_{i=0}^{n-1}F(Y_i,Y_{i+1})\right)
 \leq\frac{1+2/(1-\rho)}{n}\norm{F}_\infty^2.
\end{equation}
The analogous state-observable bound follows with $P^a$. These inequalities
apply separately to each bandwidth-dependent observable; independence of
pairs is never assumed.

\begin{lemma}[Population regularity and empirical limits]\label{empirical}
The kernel $q$ is smooth on $\overline D\times\overline D$. Almost surely,
\begin{equation}\label{strongkernellimit}
 \widetilde\mu_j\longrightarrow\mu\text{ in }C^2(\overline D),\qquad
 q_j\longrightarrow q\text{ in }C^{2,2}(\overline D\times\overline D),
 \qquad \mu_j\longrightarrow\mu\text{ in }C^2(\overline D).
\end{equation}
\end{lemma}
\begin{proof}
The smooth elliptic conormal realization has compact resolvent. If
$-Lu_k=\gamma_k u_k$, the min--max principle and \eqref{bounds} compare
$\gamma_k$ with the Neumann Laplacian eigenvalues; their counting is
polynomial. Smooth elliptic graph regularity and Sobolev embedding give
\begin{equation}\label{graph}
 \norm{f}_{H^{2a}(D)}\leq B_a\norm{(\Id-L)^a f}_{L^2(\mu)},
 \qquad
 \norm{u_k}_{C^r(\overline D)}\leq B_r(1+\gamma_k)^{b_r}
\end{equation}
for sufficiently high integer $a$ and some finite $b_r$. The graph estimates
are the standard smooth conormal elliptic estimates; see
\cite[Theorem~2.2 and Corollary~2.3]{grubb} and the Sobolev embedding
\cite[Corollary~6.12]{sobolev}. The weighted and ordinary $L^2$ norms are
equivalent; use the positive shift $\Id-L$ to include the constant mode.
Thus every fixed differentiated series
\[
 q(x,y)=\mu(x)\mu(y)\sum_k e^{-\Delta\gamma_k}u_k(x)u_k(y)
\]
converges absolutely uniformly. This proves the asserted positive-time
regularity, using population elliptic theory rather than estimated modes.

Reversibility and stationarity give
$\E\widetilde\mu_j=R_h\mu$ and
$\E q_j=(R_h\otimes R_h)q$; their deterministic biases vanish in the
displayed norms by Lemma~\ref{kernel}.
For $|\alpha|,|\beta|\leq2$, each differentiated summand of
\eqref{pair} is bounded by a constant times $(j+1)^{2d+4}$. By
\eqref{variance} its average has variance at each $(x,y)$ at most
$B(j+1)^{4d+8}2^{-j}$. A deterministic net of
$\overline D\times\overline D$ with spacing
$\varepsilon_j=2^{-j/(4d)}$ has at most $B2^{j/2}$ points.
Chebyshev's inequality and a union bound yield
\begin{equation}\label{pairprob}
 \Prob\!\left(\max_{\text{net}}|
 \partial_x^\alpha\partial_y^\beta(q_j-\E q_j)|>(j+1)^{-1}\right)
 \leq B(j+1)^{4d+10}2^{-j/2}.
\end{equation}
The bound is summable. The derivatives in question, and their expectations,
are Lipschitz on the fixed neighborhood with bound
$B(j+1)^{2d+5}$, by one further derivative in \eqref{kernelbound}.
Interpolation from the net costs at most
$B(j+1)^{2d+5}\varepsilon_j\to0$. The neighborhood makes the short line
segments legitimate even if $D$ is not convex. Borel--Cantelli gives
uniform convergence of the centered derivatives.

For \eqref{density}, the state-observable bound, derivative bound
$B(j+1)^{d+2}$ and a $d$-dimensional net with the same spacing give the
summable probability bound
$B(j+1)^{2d+6}2^{-3j/4}$ at threshold $(j+1)^{-1}$.
There are only finitely many derivative indices, so all the asserted limits
hold on one event of probability one. Uniform convergence eventually places
$\widetilde\mu_j$ in $[3c/4,3C/2]$, making saturation inactive. Then
$Z_j\to\int_D\mu\dd=1$, proving the last limit. No mixing parameter or
truth-dependent regularity constant was used in the tuning sequence.
\end{proof}

\section{Basis-invariant deterministic reconstruction}
The argument below also explains why a fixed number of modes or individual
eigenvector consistency is unnecessary. On $H=L^2(D,dx)$, let $T$ have
kernel $q(x,y)/\sqrt{\mu(x)\mu(y)}$. Multiplication by $\sqrt\mu$ conjugates
$P$ to $T$; its orthonormal eigenfunctions are $\phi_k=\sqrt\mu u_k$,
with eigenvalues $\kappa_k=e^{-\Delta\gamma_k}>0$, including the stationary
constant eigenfunction. Zero is only an accumulation point. Put
$M_j=\mu_j^{-1/2}$, $M=\mu^{-1/2}$ and $W_j=M_jT_j$, $W=MT$.

Fix $E\Subset D$. Lemma~\ref{empirical}, the common density lower bound,
and Cauchy--Schwarz in the kernel's second variable imply
\begin{equation}\label{operatorlimit}
 \norm{T_j-T}_\op\longrightarrow0,\qquad
 \norm{W_j-W}_{H\to C^2(E)}\longrightarrow0.
\end{equation}
For clarity, the needed deterministic hypotheses are merely uniform
convergence of $\mu_j$, local $C^2$ convergence of $\mu_j$, global
$L^2$ convergence of $q_j$, and
$\sup_{x\in E}\norm{\partial_x^\alpha(q_j-q)(x,\cdot)}_{L^2(D)}\to0$
for $|\alpha|\leq2$. Multiplication and its derivatives are controlled by
the positive lower bound. The global kernel bound gives Hilbert--Schmidt,
hence operator, convergence of $T_j$; differentiating the $W_j$ kernel gives
the second assertion of \eqref{operatorlimit}.

Define bounded operators $B_j(x):H\to\R^m$ and $e_j(x):H\to\R$ by
$B_j(x)f=J(W_j f)(x)$ and $e_j(x)f=(W_j f)(x)$, and similarly $B,e$.
Their norms converge uniformly on $E$. Set
\begin{equation}\label{functions}
 g(t)=\begin{cases}t^2&t>0,\\0&t\leq0,\end{cases}\qquad
 h(t)=\begin{cases}t^2\log t&t>0,\\0&t\leq0.\end{cases}
\end{equation}
Both are continuous at zero. The normal equation for \eqref{fit} is
\begin{equation}\label{normal}
 \widehat\theta_j=(A_j+\lambda_j\Id)^{-1}r_j,\qquad
 A_j=B_jg(T_j)B_j^*,\qquad
 r_j=\frac{\mu_j}{\Delta}B_jh(T_j)e_j^*.
\end{equation}
Indeed $W_j\phi_{j,k}=\kappa_{j,k}u_{j,k}$, so the outer factors supply
two eigenvalue powers and the middle function supplies two more. These
identities apply to the entire finite spectrum, including negative modes,
without taking their logarithms. They are invariant under orthogonal basis
changes within repeated eigenspaces.

The spectra lie eventually in a common bounded interval by
\eqref{operatorlimit}. Polynomial approximation of continuous functions on
that interval, and continuity of each polynomial in a bounded operator,
show that $g(T_j)\to g(T)$ and $h(T_j)\to h(T)$ in operator norm. Therefore
\begin{equation}\label{normalconvergence}
 A_j\longrightarrow A=Bg(T)B^*,\qquad
 r_j\longrightarrow r=(\mu/\Delta)Bh(T)e^*
 \quad\text{uniformly on }E.
\end{equation}
This includes empirical eigenvalue splits, crossings through zero and small
positive eigenvalues. Initial eigenvalues greater than one cause no problem
and are not interpreted as population generator eigenvalues.

\begin{proposition}[Full spectral-jet excitation]\label{excitation}
For every $x\in D$, $A(x)$ is positive definite and $r(x)=A(x)\theta(x)$.
For each compact $E\Subset D$ there is a truth-dependent $a_E>0$ with
$A(x)\geq a_E\Id$ on $E$.
\end{proposition}
\begin{proof}
Spectral expansion of the bounded-functional formulas gives
\[
 A(x)=\sum_k\kappa_k^4 Ju_k(x)Ju_k(x)^\tr,\quad
 r(x)=\sum_k\kappa_k^4\nu_k\mu(x)u_k(x)Ju_k(x).
\]
These series converge by the smoothing already proved (or by Parseval for
the bounded functionals). Equation~\eqref{pde} gives $r=A\theta$.
If $\xi^\tr A(x)\xi=0$, each nonnegative term vanishes and
$\xi^\tr Ju_k(x)=0$ for every $k$, since all $\kappa_k>0$.
Every $f\in C_c^\infty(D)$ belongs to every power domain of $\Id-L$:
successive applications of the local differential operator retain interior
support and satisfy all iterated conormal conditions. Its spectral partial
sums converge in each graph norm. Equation~\eqref{graph} and Sobolev
embedding, choosing $2a>d/2+2$, imply convergence in $C^2(\overline D)$.
Thus $\xi^\tr Jf(x)=0$ for every such $f$.

Multiply affine and quadratic polynomials centered at $x$ by a cutoff equal
to one near $x$. Their gradients and symmetric Hessians can be prescribed
arbitrarily; the resulting jets span $\R^m$. Hence $\xi=0$.
Finally $A(x)$ is continuous by continuity of the evaluation functionals;
compactness of $E$ gives $a_E>0$. This proves excitation rather than assuming
an unverified rank condition. No class-uniform conditioning bound is claimed.
\end{proof}

\begin{proof}[Consistency in Theorem~\ref{main}]
Work on the event of Lemma~\ref{empirical}. By
\eqref{normalconvergence}, eventually $A_j\geq a_E\Id/2$ uniformly on
$E$. The exact identity
\[
 \widehat\theta_j-\theta=(A_j+\lambda_j\Id)^{-1}
 \big[(r_j-r)+(A-A_j)\theta-\lambda_j\theta\big]
\]
proves local uniform convergence for any positive $\lambda_j\to0$.
No relative noise/ridge rate is needed because the limiting local matrix is
positive definite. This also proves the stated local drift and covariance
limits, using the density lower bound.

Projection onto the closed convex matrix interval in Frobenius norm is
nonexpansive and fixes $S$. On a countable compact exhaustion of $D$, the
projected tensor converges pointwise in the interior and is uniformly bounded
by the known matrix interval. Since $D$ has finite volume, dominated
convergence gives the global $L^2$ limit. There is no corresponding global
claim about the fitted divergence or derivatives of the tensor estimate.
All sample sizes use a dyadic index tending to infinity. Measurability is
proved next, completing the theorem.
\end{proof}

\section{Finite-feature measurability and verification}
The eigenbasis in \eqref{empiricalmodes} need not be chosen measurably. The
following equivalent formulas avoid it. For $0\leq i\leq n$ put
$f_i=K_h(\cdot,Y_i)/\sqrt{\mu_j}$ and
$v_i=K_h(\cdot,Y_i)/\mu_j$. Let $V:\R^{n+1}\to H$ have columns $f_i$,
and let $C_n$ be the symmetric matrix with adjacent entries
$(C_n)_{i,i+1}=(C_n)_{i+1,i}=1/(2n)$ and all others zero. Then
\[
 T_j=VC_nV^*,\qquad G=V^*V\geq0,\qquad H_n=G^{1/2}C_nG^{1/2}.
\]
Let $D_x$ have columns $Jv_i(x)$ and let $v(x)$ be the column of $v_i(x)$.
The formulas in \eqref{normal} become
\begin{align}
 A_j(x)&=D_xC_nG^{1/2}g(H_n)G^{1/2}C_nD_x^\tr,\label{gram}\\
 r_j(x)&=\frac{\mu_j(x)}{\Delta}
 D_xC_nG^{1/2}h(H_n)G^{1/2}C_nv(x).\notag
\end{align}
To check them, $B_j=D_xC_nV^*$ and $e_j=v(x)^\tr C_nV^*$, and
\[
 V^*a(VC_nV^*)V=G^{1/2}a(H_n)G^{1/2}
\]
holds first for polynomials $a$ and then for continuous $a$ by uniform
approximation. No inverse of $G$ is used, so singular feature Gram matrices
are allowed. Kernel sections and their derivatives are Borel functions of
the data; known-domain integrals defining $Z_j$ and $G$ are Borel as well.
Positive square root and continuous functional calculus of finite symmetric
matrices are continuous, and $A_j+\lambda_j\Id$ is invertible.
Thus \eqref{gram} proves joint measurability and continuity in $x$.
Separability on compact sets, and the bounded projected output, give the
usual $C(E)$ and $L^2(D)$ random-element measurability. Multiplicities,
rank loss and zero crossings require no eigenvector selection.

The supplement contains the source, a portable finite-control runner,
expected scientific outputs, exact source-edition notes, and the separately
verified classical regularization application. Its controls check algebra,
signed-kernel moments, covariance identities, spectral collisions and
boundary/model distinctions. They supplement the analytic proof and do not
prove an infinite-dimensional theorem, simulate the estimator's asymptotic
performance, or supply a general-domain numerical implementation.

\paragraph{AI use and review status.}
AI tools were used extensively in deriving, checking and writing this work,
including separate adversarial mathematical and priority audits. Their
checks are disclosed as AI-assisted verification, not independent human
peer review or a proof-assistant certificate. This preprint is unrefereed
and has not undergone human peer review. The author takes responsibility for
the claims and for the explicit scope and source-access limitations.

\begin{thebibliography}{99}
\bibitem{reiss} M.~Rei\ss{} (with J.~Chorowski, E.~Gobet and M.~Hoffmann),
SDE estimation from discrete observations as inverse problems, in
\emph{Oberwolfach Reports} 24/2017, 1507--1509. Report published April 28,
2018. \href{https://doi.org/10.4171/OWR/2017/24}{doi:10.4171/OWR/2017/24}.
\bibitem{hs} L.~P.~Hansen and J.~A.~Scheinkman,
\emph{Back to the Future: Generating Moment Implications for Continuous-Time
Markov Processes}, NBER Technical Working Paper 141, September 1993.
Inspected working-paper Proposition~5.5.
\href{https://www.nber.org/system/files/working_papers/t0141/t0141.pdf}{Primary working paper}.
\bibitem{chs} X.~Chen, L.~P.~Hansen and J.~Scheinkman,
Nonlinear principal components and long-run implications of multivariate
diffusions, \emph{Ann. Statist.} 37 (2009), 4279--4312.
\href{https://doi.org/10.1214/09-AOS706}{doi:10.1214/09-AOS706}.
\bibitem{cv2006} D.~Crommelin and E.~Vanden-Eijnden,
Reconstruction of diffusions using spectral data from timeseries,
\emph{Commun. Math. Sci.} 4 (2006), 651--668.
\href{https://doi.org/10.4310/CMS.2006.v4.n3.a9}{doi:10.4310/CMS.2006.v4.n3.a9}.
\bibitem{cv2011} D.~Crommelin and E.~Vanden-Eijnden,
Diffusion estimation from multiscale data by operator eigenpairs,
\emph{Multiscale Model. Simul.} 9 (2011), 1588--1623.
\href{https://doi.org/10.1137/100795917}{doi:10.1137/100795917}.
Body inspected: \href{https://ir.cwi.nl/pub/18934/18934B.pdf}{CWI author
manuscript}, revised August 31, 2011; publisher-final body uninspected.
\bibitem{ghr} E.~Gobet, M.~Hoffmann and M.~Rei\ss{},
Nonparametric estimation of scalar diffusions based on low frequency data,
\emph{Ann. Statist.} 32 (2004), 2223--2253.
\href{https://doi.org/10.1214/009053604000000797}{doi:10.1214/009053604000000797}.
\bibitem{chorowski} J.~Chorowski,
Nonparametric volatility estimation in scalar diffusions: Optimality across
observation frequencies, \emph{Bernoulli} 24 (2018), 2934--2990.
\href{https://doi.org/10.3150/17-BEJ950}{doi:10.3150/17-BEJ950}.
Inspected \href{https://arxiv.org/abs/1507.07139v2}{arXiv:1507.07139v2},
March 31, 2016.
\bibitem{nickl} R.~Nickl,
Consistent inference for diffusions from low frequency measurements,
\emph{Ann. Statist.} 52 (2024), 519--549.
\href{https://doi.org/10.1214/24-AOS2357}{doi:10.1214/24-AOS2357}.
\bibitem{gw} M.~Giordano and S.~Wang,
Statistical algorithms for low-frequency diffusion data: A PDE approach,
\emph{Ann. Statist.} 53 (2025), 1150--1175.
\href{https://doi.org/10.1214/25-AOS2496}{doi:10.1214/25-AOS2496}.
\bibitem{abjn} G.~S.~Alberti, D.~Barnes, A.~Jambhale and R.~Nickl,
On low frequency inference for diffusions without the hot spots conjecture,
\emph{Math. Stat. Learn.} 8 (2025), 305--322.
\href{https://doi.org/10.4171/MSL/53}{doi:10.4171/MSL/53}.
Inspected \href{https://arxiv.org/abs/2410.19393v2}{arXiv:2410.19393v2},
July 22, 2025, especially Lemma~3.8.
\bibitem{bal} G.~Bal and G.~Uhlmann,
Reconstruction of coefficients in scalar second-order elliptic equations
from knowledge of their solutions, \emph{Comm. Pure Appl. Math.} 66 (2013),
1629--1652. \href{https://doi.org/10.1002/cpa.21453}{doi:10.1002/cpa.21453}.
Inspected \href{https://www.stat.uchicago.edu/~guillaumebal/PAPERS/Hybrid-Meas-u-12.pdf}{author manuscript},
April 10, 2013; publisher-final identity not assumed.
\bibitem{kohn} R.~V.~Kohn and B.~D.~Lowe,
A variational method for parameter identification,
\emph{M2AN} 22 (1988), 119--158, especially Section~5.
\href{https://www.numdam.org/item/M2AN_1988__22_1_119_0/}{Primary journal archive}.
\bibitem{grubb} G.~Grubb,
Regularity of spectral fractional Dirichlet and Neumann problems,
\emph{Math. Nachr.} 289 (2016), 831--844.
\href{https://doi.org/10.1002/mana.201500041}{doi:10.1002/mana.201500041}.
\bibitem{sobolev} G.~Grubb, \emph{Distributions and Operators},
Chapter~6, Sobolev spaces, Corollary~6.12.
\href{https://web.math.ku.dk/~grubb/dist6.pdf}{Author's chapter}.
\bibitem{lenzen} F.~Lenzen and O.~Scherzer,
Tikhonov type regularization methods: History and recent progress,
\emph{ECCOMAS 2004}, especially Section~3, Theorem~3.2.
\href{https://franklenzen.de/pdf/lenzen_scherzer_eccomas2004.pdf}{Author's full paper}.
\end{thebibliography}
\end{document}
